\section{Low-Resource Training System}

\subsection{Base model}
The base model is Qwen3.5-2B, a vision-language model released by the Qwen team under an Apache-2.0 license \citep{qwen3_5}. It pairs a text transformer with a vision encoder through early-fusion multimodal training, and its language stack combines Gated DeltaNet layers, gated attention, and a sparse mixture-of-experts feed-forward network, with a native context of 262{,}144 tokens. The checkpoint contains 1{,}892{,}736{,}832 parameters in total.

We adapt only the text transformer. The vision encoder is left untouched and is never exercised: every dataset in this study is text-only, and the model is prompted purely as a language model. The description that follows therefore concerns the language stack alone.

\subsection{Adaptation}
We fine-tune with QLoRA \citep{dettmers2023qlora}, which keeps the base weights frozen and quantized to 4-bit NF4 while training low-rank adapters on top \citep{hu2021lora}. Adapters of rank 16 with scaling factor $\alpha = 32$ and no dropout are attached to the query, key, value, output, gate, up, and down projections of every transformer block. This leaves 10{,}911{,}744 parameters trainable, 0.58\% of the model. Attention uses a scaled dot-product implementation. Table~\ref{tab:trainconfig} lists the configuration.

\begin{table}[t]
\centering
\caption{Training configuration. The corpus and its provenance are described in Section~\ref{sec:data}.}
\label{tab:trainconfig}
\begin{tabular}{ll}
\toprule
Setting & Value \\
\midrule
Base model & Qwen3.5-2B (text transformer) \\
Total parameters & 1{,}892{,}736{,}832 \\
Trainable parameters & 10{,}911{,}744 (0.58\%) \\
Base quantization & 4-bit NF4 \\
LoRA rank / $\alpha$ / dropout & 16 / 32 / 0 \\
Target modules & q, k, v, o, gate, up, down projections \\
Attention & scaled dot-product \\
Sequence length & 384 \\
Batch size $\times$ grad.\ accumulation & $1 \times 8$ \\
Epochs / optimizer steps & 1 / 2{,}593 \\
Learning rate & $1.5\times10^{-4}$ \\
Training examples & 20{,}741 \\
\bottomrule
\end{tabular}
\end{table}

\subsection{Fitting the run into 4~GB}
Several choices keep the run inside the memory of a single RTX A2000 Laptop GPU with 4~GB of VRAM, on Windows~11. The base is quantized to 4-bit, so the frozen weights occupy roughly 1~GB rather than the four they would take in half precision. A short sequence length of 384 tokens caps the activation and attention-cache footprint, which is the term that grows with context. A batch size of one with gradient accumulation over eight steps gives an effective batch of eight without holding eight sequences in memory at once. Periodic in-loop evaluation is disabled: the model has a large vocabulary, and materializing logits over it for an evaluation batch overflows the remaining memory, so we hold out 200 examples but score them after training rather than during it.

Under these settings the run peaks at 1.81~GB of VRAM, which leaves headroom on a 4~GB card and keeps inference on the same hardware comfortable. Training a single epoch of 2{,}593 steps takes 100.5 minutes and ends at a training loss of 0.4145. The resulting adapter is 41.6~MiB on disk.
